\chapter{Verilerin Grafiksel ve Sayısal Betimlenmesi}
\label{ch:descriptive}
\index{betimsel istatistik}

\begin{hedefler}
Bu bölümde frekans dağılımları, temel grafikler, merkezi eğilim ve yayılım
ölçüleri birlikte ele alınır. Amaç yalnız sayı hesaplamak değil, dağılımın
merkezini, değişkenliğini, biçimini ve olağandışı gözlemlerini açıklamaktır.
\end{hedefler}

\begin{hazirlik}
$(2,3,3,4,18)$ veri dizisinde ortalama ile medyandan hangisinin tipik değeri
daha iyi temsil edeceğini hesaplamadan önce tahmin edin ve gerekçelendirin.
\end{hazirlik}

\section{PDR vakası: Akademik kaygı verisini ilk kez görmek}

Bir psikolojik danışma birimi, sınav döneminde öğrencilerin akademik kaygı
puanlarını incelemek istemektedir. Araştırmacının ilk sorusu ``Hangi test
kullanılmalı?'' değildir. Önce verinin kimlerden ve nasıl toplandığı, hangi
değişkenlerin bulunduğu, eksik veya olanaksız değer olup olmadığı ve puanların
nasıl dağıldığı anlaşılmalıdır.

Bu bölüm boyunca aşağıdaki sorulara yanıt aranacaktır:

\begin{enumerate}
  \item Kaygı puanlarının tipik düzeyi nedir?
  \item Öğrenciler birbirine ne kadar benzemekte veya birbirinden ne kadar ayrılmaktadır?
  \item Dağılım simetrik midir; aykırı değer adayı var mıdır?
  \item Hangi grafik ve sayısal özet araştırma sorusuna en açık yanıtı verir?
\end{enumerate}

Betimsel analiz yalnız çıkarımsal analize hazırlık değildir. Örneklemde ne
gözlendiğini dürüst ve anlaşılır biçimde açıklayan başlı başına bir araştırma
sonucudur.

\section{Önce veri yapısını tanımak}
\index{veri yapısı}

Kategorik değişkenler frekans ve yüzdelerle; nicel değişkenler ise sıralı veri,
grafikler ve sayısal özetlerle incelenir. Bir grafik seçilirken değişkenin
ölçme düzeyi belirleyicidir. Aynı sayısal özetlerin çok farklı veri
örüntülerini gizleyebilmesi nedeniyle grafiksel ve sayısal inceleme birlikte
yürütülmelidir \cite{anscombe1973}.

\subsection{Analiz birimi ve kod kitabı}
\index{analiz birimi}
\index{kod kitabı}

Her satırın kimi veya neyi temsil ettiği \emph{analiz birimi} olarak açıkça
belirlenmelidir. Bir satır öğrenci, görüşme, okul ya da ölçüm zamanı olabilir.
Aynı öğrencinin ön-test ve son-test ölçümleri iki bağımsız öğrenci gibi
yorumlanamaz.

Analize başlamadan önce kısa bir kod kitabı hazırlanması yararlıdır:

\begin{table}[htbp]
\centering
\caption{Örnek PDR veri dosyası için kod kitabı.}
\begin{tabular}{@{}p{0.20\textwidth}p{0.24\textwidth}p{0.17\textwidth}p{0.23\textwidth}@{}}
\toprule
Değişken & Açıklama & Tür/düzey & Olası değerler \\
\midrule
\texttt{id} & Katılımcı kodu & Kimlik & 1, 2, $\ldots$, 120 \\
\texttt{sinif} & Sınıf düzeyi & Sıralı kategorik & 1, 2, 3, 4 \\
\texttt{kaygi} & Akademik kaygı toplam puanı & Nicel & 20--100 \\
\texttt{basvuru} & Danışmanlığa başvuru & İkili kategorik & 0: Hayır, 1: Evet \\
\bottomrule
\end{tabular}
\end{table}

Kod kitabındaki olası değer aralığı veri temizliğini de yönlendirir. Kaygı
puanı için 140 görülmesi, doğrudan aykırı değer analizi yapılmasından önce bir
veri giriş veya puanlama hatası olasılığını düşündürmelidir.

\section{Frekans, oran ve yüzde}
\index{frekans}
\index{oran}
\index{yüzde}
\index{frekans tablosu}

Bir kategorinin frekansı $f_j$, o kategoriye düşen gözlem sayısıdır. Göreli
frekans ve yüzde
\[
r_j=\frac{f_j}{n},\qquad \%_j=100\frac{f_j}{n}
\]
biçiminde hesaplanır. Örneğin 80 öğrencinin 28'i danışmanlık hizmetine
başvurmuşsa başvuru yüzdesi $100(28/80)=\%35$ olur.

Paydanın açık olması önemlidir. Beş öğrencinin başvuru bilgisi eksikse yüzde
toplam 80 öğrenciye göre mi, yoksa geçerli 75 yanıta göre mi hesaplanmıştır?
Rapor bunu belirtmelidir. Frekans tablosunda sayı ve yüzdeyi birlikte vermek,
özellikle grupların büyüklükleri farklı olduğunda yanlış yorum riskini azaltır.

\begin{mistakebox}
Yüzdeleri küçük örneklemlerde bağlamdan koparmayın. Dört kişilik bir grupta
tek katılımcı yüzde 25'e karşılık gelir. ``Katılımcıların yüzde 25'i'' ifadesi,
``dört katılımcıdan biri'' bilgisiyle birlikte verilmelidir.
\end{mistakebox}

\begin{table}[htbp]
\centering
\caption{Değişken türüne göre grafik seçimi.}
\begin{tabular}{p{0.27\textwidth}p{0.57\textwidth}}
\toprule
Veri yapısı & Önerilen grafik \\
\midrule
Tek kategorik değişken & Çubuk grafiği, Pareto grafiği \\
Tek nicel değişken & Histogram, nokta grafiği, kutu grafiği \\
İki nicel değişken & Saçılım grafiği \\
Nicel sonuç ve kategorik grup & Yan yana kutu grafikleri \\
İki kategorik değişken & Çapraz tablo, yığılmış çubuk veya mozaik grafiği \\
\bottomrule
\end{tabular}
\end{table}

\begin{etkinlik}
Bir PDR araştırmasında şu değişkenler ölçülsün: cinsiyet, sınıf düzeyi, kaygı
puanı, haftalık uyku süresi ve danışmanlığa başvuru durumu. Her değişken için
uygun bir grafik seçin. Seçiminizi ölçme düzeyi ve araştırma sorusuyla
gerekçelendirin.
\end{etkinlik}

\begin{figure}[htbp]
\centering
\begin{tikzpicture}[alt={Yaklaşık simetrik bir puan dağılımını gösteren histogram},x=0.72cm,y=0.55cm]
  \draw[->,PrismGray] (0,0) -- (8.7,0) node[right,font=\scriptsize]{puan};
  \draw[->,PrismGray] (0,0) -- (0,5.2) node[above,font=\scriptsize]{frekans};
  \foreach \x/\h in {0.4/1.0,1.2/2.1,2.0/3.8,2.8/4.7,3.6/4.0,4.4/2.7,5.2/1.5,6.0/0.8}
    \filldraw[fill=PrismSubheadingBlue!42,draw=PrismBlue] (\x,0) rectangle +(0.72,\h);
  \draw[PrismWarnOrange!85!black,very thick] (7.7,0) -- (7.7,1.0);
  \fill[PrismWarnOrange!85!black] (7.7,1.0) circle (2pt);
  \node[font=\scriptsize,align=center,text=PrismGray] at (7.55,1.65) {aykırı değer\\adayı};
  \draw[PrismBlue,dashed] (3.15,0) -- (3.15,4.95);
  \node[font=\scriptsize,text=PrismBlue] at (3.15,5.25) {merkez};
\end{tikzpicture}
\caption{Dağılımı merkez, yayılım, biçim ve aykırı değerlerle birlikte okuma.}
\end{figure}

\section{Merkezi eğilim ölçüleri}
\index{merkezi eğilim}
\index{ortalama}
\index{medyan}
\index{mod}

$x_1,\ldots,x_n$ gözlemleri için aritmetik ortalama
\[
\bar{x}=\frac{1}{n}\sum_{i=1}^{n}x_i
\]
olarak tanımlanır. Medyan, sıralanmış veriyi iki eşit parçaya ayıran değerdir.
Mod en sık gözlenen değer veya kategoridir.

Medyan hesaplanırken veriler önce küçükten büyüğe sıralanır. $n$ tek ise
ortadaki değer, $n$ çift ise ortadaki iki değerin ortalaması alınır. Modun
birden fazla olması mümkündür; hiçbir değer tekrar etmiyorsa belirgin bir mod
bulunmayabilir.

\subsection{Hangi merkez ölçüsü?}

\begin{description}
  \item[Ortalama] Nicel, yaklaşık simetrik ve ciddi aykırı değeri olmayan veride etkilidir.
  \item[Medyan] Çarpık dağılımda, aykırı değer varlığında veya sıralı veride daha sağlamdır.
  \item[Mod] En yaygın kategori ya da puanın araştırma açısından anlamlı olduğu durumda kullanılır.
\end{description}

Ortalama bütün gözlemleri kullandığı için aykırı değerlerden etkilenir. Medyan
ise sıralamaya dayandığı için daha sağlamdır. Sağa çarpık bir bekleme süresi
dağılımında medyan, tipik katılımcının deneyimini ortalamadan daha iyi
yansıtabilir.

\section{Yayılım ölçüleri}
\index{yayılım}
\index{değişim aralığı}
\index{varyans}
\index{standart sapma}
\index{çeyrekler arası aralık}

Örneklem varyansı ve standart sapması
\[
s^2=\frac{1}{n-1}\sum_{i=1}^{n}(x_i-\bar{x})^2,
\qquad s=\sqrt{s^2}
\]
biçimindedir. Açıklık $\max(x)-\min(x)$ ile, çeyrekler arası açıklık ise
$IQR=Q_3-Q_1$ ile hesaplanır.

Standart sapma ortalamayla birlikte, IQR ise çoğunlukla medyanla birlikte
raporlanır. Dağılım yaklaşık simetrik ve aykırı değerlerden uzaksa
$\bar{x}\pm s$; belirgin biçimde çarpıksa medyan ve IQR daha bilgilendirici bir
ikili oluşturur.

\subsection{Yayılım neden merkez kadar önemlidir?}

İki danışmanlık grubunun kaygı ortalaması 50 olabilir. Birinci grupta puanlar
48--52 arasında toplanırken ikinci grupta 25--75 arasında yayılabilir. Yalnız
ortalamayı raporlamak bu önemli farkı gizler.

Standart sapma gözlemlerin ortalamadan tipik uzaklığını aynı ölçü biriminde
özetler. Varyans hesaplama açısından yararlı olsa da puanın karesi biriminde
olduğu için doğrudan yorumlanması daha güçtür. Açıklık yalnız iki uç değere
dayanır; IQR ise ortadaki yüzde 50'nin yayıldığı alanı gösterdiğinden uç
değerlere karşı daha sağlamdır.

\begin{sezgibox}
Merkez ``veri nerede toplanıyor?'', yayılım ise ``gözlemler birbirinden ne
kadar farklı?'' sorusunu yanıtlar. Aynı ortalamaya sahip iki grup çok farklı
standart sapmalara sahip olabilir.
\end{sezgibox}

\section{Sınıflandırılmış verilerde yer ve yayılma ölçüleri}
\index{sınıflandırılmış veri}
\index{gruplandırılmış veri}
\index{sınıf orta noktası}
\index{kümülatif frekans}

Gözlem sayısı fazla olduğunda nicel değerler sınıf aralıklarında özetlenebilir.
Bu durumda ham puanların tek tek değerleri değil, her aralığa düşen gözlem
sayısı bilinir. Hesaplamalarda sınıftaki bütün gözlemlerin sınıf orta
noktasında bulunduğu varsayıldığı için elde edilen yer ve yayılma ölçüleri
\emph{yaklaşıktır}. Ham veri mevcutsa temel analiz ham değerlerden yapılmalıdır.

$j$'inci sınıfın alt ve üst sınırları $a_j$ ve $b_j$, frekansı $f_j$ olsun.
Sınıf orta noktası ve toplam gözlem sayısı
\[
m_j=\frac{a_j+b_j}{2},
\qquad
n=\sum_{j=1}^{k}f_j
\]
ile bulunur. Kesintisiz sınıflarda aralıkların boşluk bırakmaması gerekir.
Tam sayı puanlarda 20--29 ve 30--39 sınıflarının gerçek sınırları sırasıyla
19,5--29,5 ve 29,5--39,5 olarak düşünülebilir.

\subsection{Sınıflandırılmış ortalama, varyans ve standart sapma}

Sınıflandırılmış verinin yaklaşık ortalaması ağırlıklı ortalamadır:
\[
\bar{x}_g=\frac{\sum_{j=1}^{k}f_jm_j}{n}.
\]
Yaklaşık örneklem varyansı ve standart sapması
\[
s_g^2=\frac{\sum_{j=1}^{k}f_j(m_j-\bar{x}_g)^2}{n-1}
=\frac{\sum f_jm_j^2-n\bar{x}_g^2}{n-1},
\qquad
s_g=\sqrt{s_g^2}
\]
biçiminde hesaplanır. Frekansı yüksek sınıflar ortalama ve varyans hesabına
daha fazla ağırlık verir.

\subsection{Sınıflandırılmış medyan, çeyrekler ve mod}

Önce kümülatif frekanslar hesaplanır. $n/2$'inci gözlemi içeren aralık
\emph{medyan sınıfıdır}. Sınıf içinde gözlemlerin yaklaşık eşit dağıldığı
varsayımıyla
\[
Md_g=L+\left(\frac{n/2-F_{\mathrm{önce}}}{f_{Md}}\right)h
\]
kullanılır. Burada $L$ medyan sınıfının gerçek alt sınırı,
$F_{\mathrm{önce}}$ önceki sınıfın kümülatif frekansı, $f_{Md}$ medyan
sınıfının frekansı ve $h$ sınıf genişliğidir. Aynı formülde $n/2$ yerine
$n/4$ ve $3n/4$ kullanılarak $Q_1$ ve $Q_3$ yaklaşık bulunur.

En yüksek frekanslı aralık mod sınıfıdır. Komşu sınıflar dikkate alınarak
yaklaşık mod
\[
Mo_g=L+\frac{d_1}{d_1+d_2}h,
\qquad
d_1=f_m-f_{m-1},\quad d_2=f_m-f_{m+1}
\]
ile hesaplanabilir. Açıklık için en üst gerçek sınıf sınırından en alt gerçek
sınıf sınırı çıkarılabilir; ancak ilk veya son sınıf açık uçluysa açıklık ve
orta nokta temelli ölçüler hesaplanamaz.

\subsection{Adım adım çözümlü sınıflandırılmış veri örneği}

Elli öğrencinin akademik kaygı puanları aşağıdaki frekans dağılımıyla
özetlenmiş olsun:

\begin{table}[htbp]
\centering
\caption{Akademik kaygı puanlarının sınıflandırılmış frekans dağılımı.}
\begin{tabular}{@{}lrrrrr@{}}
\toprule
Puan sınıfı & Gerçek sınırlar & $m_j$ & $f_j$ & $F_j$ & $f_jm_j$ \\
\midrule
20--29 & 19,5--29,5 & 24,5 & 3  & 3  & 73,5 \\
30--39 & 29,5--39,5 & 34,5 & 7  & 10 & 241,5 \\
40--49 & 39,5--49,5 & 44,5 & 12 & 22 & 534,0 \\
50--59 & 49,5--59,5 & 54,5 & 15 & 37 & 817,5 \\
60--69 & 59,5--69,5 & 64,5 & 9  & 46 & 580,5 \\
70--79 & 69,5--79,5 & 74,5 & 4  & 50 & 298,0 \\
\midrule
Toplam & & & 50 & & 2545,0 \\
\bottomrule
\end{tabular}
\end{table}

\adimbaslik{1. Yaklaşık ortalama}

\[
\bar{x}_g=\frac{2545}{50}=50{,}90.
\]

\adimbaslik{2. Yaklaşık varyans ve standart sapma}

Orta noktaların frekanslarla ağırlıklandırılmış kareli sapmaları toplamı
$8552$'dir. Buna göre
\[
s_g^2=\frac{8552}{49}\approx174{,}53,
\qquad
s_g\approx13{,}21.
\]

\adimbaslik{3. Yaklaşık medyan ve çeyrekler}

$n/2=25$'inci gözlem, kümülatif frekansı 22'den 37'ye çıkaran 50--59
sınıfındadır. Dolayısıyla
\[
Md_g=49{,}5+\left(\frac{25-22}{15}\right)10=51{,}50.
\]
$n/4=12{,}5$'inci gözlem 40--49 sınıfında,
$3n/4=37{,}5$'inci gözlem 60--69 sınıfındadır:
\[
Q_{1g}\approx39{,}5+\left(\frac{12{,}5-10}{12}\right)10=41{,}58,
\]
\[
Q_{3g}\approx59{,}5+\left(\frac{37{,}5-37}{9}\right)10=60{,}06.
\]
Böylece $IQR_g\approx60{,}06-41{,}58=18{,}48$ olur.

\adimbaslik{4. Yaklaşık mod ve açıklık}

Mod sınıfı 50--59'dur. $d_1=15-12=3$ ve $d_2=15-9=6$ olduğundan
\[
Mo_g=49{,}5+\frac{3}{3+6}(10)\approx52{,}83.
\]
Gerçek sınıf sınırlarına dayalı yaklaşık açıklık $79{,}5-19{,}5=60$'tır.

\begin{mistakebox}
Sınıflandırılmış ortalama, standart sapma ve yüzdelikler ham veriden elde
edilmiş kesin değerler gibi raporlanmamalıdır. Sınıf aralıkları genişledikçe
orta nokta ve sınıf içi doğrusal ara değerleme varsayımlarından kaynaklanan
bilgi kaybı büyür.
\end{mistakebox}

\subsection{Python ve SPSS ile hesaplama}

Orta noktaların frekans kadar tekrarlandığı varsayımıyla yaklaşık özetler
Python'da şöyle hesaplanabilir:

\begin{lstlisting}
import numpy as np

orta = np.array([24.5, 34.5, 44.5, 54.5, 64.5, 74.5])
frekans = np.array([3, 7, 12, 15, 9, 4])
n = frekans.sum()

ortalama = np.average(orta, weights=frekans)
varyans = np.sum(frekans * (orta - ortalama)**2) / (n - 1)
standart_sapma = np.sqrt(varyans)

print(ortalama, varyans, standart_sapma)
\end{lstlisting}

SPSS'te bir sütuna sınıf orta noktaları, ikinci sütuna frekanslar girilir;
\textbf{Data $>$ Weight Cases} menüsünde frekans değişkeniyle ağırlıklandırma
yapıldıktan sonra \textbf{Analyze $>$ Descriptive Statistics $>$ Descriptives}
kullanılabilir. Bu işlem yaklaşık ortalama ve standart sapmayı verir. Medyan
ve çeyreklerde sınıf içi ara değerleme gerektiğinden yukarıdaki formüller veya
ham veri kullanılmalıdır. İşlem tamamlandığında sonraki analizleri yanlışlıkla
ağırlıklandırmamak için \textbf{Weight Cases} seçeneği kapatılmalıdır.

\section{Dağılım biçimi ve aykırı değerler}
\index{dağılım biçimi}
\index{çarpıklık}
\index{aykırı değer}
\index{kutu grafiği}
\index{histogram}

Simetrik dağılımda iki kuyruk yaklaşık dengelidir. Sağa çarpık dağılımda uzun
kuyruk büyük değerlere, sola çarpık dağılımda küçük değerlere uzanır. Kutu
grafiğinde
\[
x<Q_1-1.5IQR \quad\text{veya}\quad x>Q_3+1.5IQR
\]
koşulunu sağlayan gözlemler aykırı değer adayı olarak işaretlenebilir.

Kutu grafiği kuralı mekanik bir silme kuralı değil, inceleme sinyalidir.
Araştırmacı aykırı değer adayının özgün forma, ölçüm zamanına ve katılımcı
bağlamına dönmelidir. Analiz sonucu bu gözleme duyarlıysa, gerekçeli biçimde
hem gözlem dahil hem hariç sonuçların gösterildiği bir duyarlılık analizi
yapılabilir.

\begin{mistakebox}
Aykırı değer etiketi, gözlemin yanlış olduğu anlamına gelmez. Veri giriş
hatası, farklı bir alt grup veya gerçek fakat seyrek bir durum söz konusu
olabilir. İnceleme yapılmadan gözlem silinmemelidir.
\end{mistakebox}

\section{Adım adım çözümlü PDR örneği}

On iki öğrencinin akademik kaygı puanları küçükten büyüğe
\[
42,48,50,51,52,53,54,55,57,60,62,85
\]
olarak sıralanmıştır.

\subsection*{1. Merkez}

Toplam $669$ olduğundan
\[
\bar{x}=\frac{669}{12}=55{,}75.
\]
Ortadaki iki değer 53 ve 54 olduğu için medyan
\[
\widetilde{x}=\frac{53+54}{2}=53{,}5
\]
olur. Ortalama medyandan büyüktür; büyük uçtaki 85 puanı ortalamayı yukarı
çekmektedir.

\subsection*{2. Yayılım}

Örneklem standart sapması $s\approx10{,}64$, açıklık $85-42=43$'tür.
Veriyi alt ve üst altışarlı yarıya ayıran yöntemle $Q_1=50{,}5$,
$Q_3=58{,}5$ ve $IQR=8$ bulunur. Üst aykırı değer sınırı
\[
Q_3+1{,}5IQR=58{,}5+1{,}5(8)=70{,}5
\]
olduğundan 85 puanı aykırı değer adayıdır.

\subsection*{3. Duyarlılık analizi}

85 puanı çıkarıldığında ortalama yaklaşık $53{,}09$, medyan 53 ve standart
sapma yaklaşık $5{,}58$ olur. Medyan çok az değişirken ortalama ve standart
sapma belirgin biçimde değişmiştir. Bu karşılaştırma 85'in mutlaka silinmesi
gerektiğini değil, sonucun bu gözleme duyarlı olduğunu gösterir.

\begin{researchbox}
PDR açısından 85 puanı yalnız istatistiksel bir sorun değildir. Puan gerçekse
yüksek kaygı yaşayan bir öğrenciyi temsil ediyor olabilir. Gözlemi sessizce
silmek hem bilimsel sonucu hem de mesleki açıdan önemli bir durumu görünmez
kılabilir.
\end{researchbox}

\section{Python ile yeniden üretilebilir analiz}

\begin{lstlisting}
import pandas as pd
import matplotlib.pyplot as plt

kaygi = pd.Series([42, 48, 50, 51, 52, 53,
                   54, 55, 57, 60, 62, 85], name="kaygi")

print(kaygi.describe())
print("Medyan:", kaygi.median())
print("IQR:", kaygi.quantile(.75) - kaygi.quantile(.25))

fig, ax = plt.subplots(1, 2, figsize=(8, 3))
kaygi.plot.hist(bins=6, edgecolor="black", ax=ax[0])
kaygi.plot.box(vert=False, ax=ax[1])
plt.tight_layout()
plt.show()
\end{lstlisting}

Yazılımın çeyrek hesaplama yöntemi elle kullanılan yöntemden küçük farklar
gösterebilir. Bu bir hata olmak zorunda değildir; farklı yazılımlar yüzdelik
değerler arasında farklı ara değerleme kuralları kullanabilir. Raporlanacak
yöntem analiz boyunca tutarlı olmalıdır.

\section{SPSS ile uygulama adımları}

\begin{enumerate}
  \item \textbf{Variable View} bölümünde \texttt{kaygi} değişkenini sayısal ve ölçek düzeyinde tanımlayın.
  \item \textbf{Analyze $>$ Descriptive Statistics $>$ Explore} yolunu izleyin.
  \item \texttt{kaygi} değişkenini \textbf{Dependent List} alanına taşıyın.
  \item \textbf{Statistics} altında betimsel istatistikleri; \textbf{Plots} altında histogram ve kutu grafiğini seçin.
  \item Çıktıda önce geçerli ve eksik $n$ değerlerini, sonra merkez, yayılım ve grafikleri inceleyin.
\end{enumerate}

\begin{outputguidebox}
Bir betimsel istatistik çıktısını şu sırayla okuyun: (1) geçerli ve eksik
gözlem sayısı, (2) minimum--maksimum ve olanaksız değerler, (3) ortalama ile
medyanın yakınlığı, (4) standart sapma ve IQR, (5) histogram ile kutu
grafiğinin sayısal özetlerle tutarlılığı. Yazılım tablosundaki tek bir sütuna
bakarak dağılım hakkında karar vermeyin.
\end{outputguidebox}

\section{Yanıltıcı grafiklerden kaçınma}
\index{yanıltıcı grafik}
\index{grafik!eksen ölçeği}

Grafiğin dikey eksenini sıfırdan başlatmamak küçük farkları olduğundan büyük
gösterebilir. Üç boyutlu sütunlar alan ve derinlik algısını bozabilir. Eşit
olmayan histogram sınıf aralıklarında yalnız sütun yüksekliğini karşılaştırmak
yanlıştır; alanın frekansı temsil etmesi gerekir. Grafik başlığı, eksen adları,
ölçü birimi, örneklem büyüklüğü ve veri kaynağı açıkça belirtilmelidir.

\begin{etkinlik}
Aynı iki grup ortalamasını gösteren iki çubuk grafiği hazırlayın: ilkinde
dikey eksen 0'dan, ikincisinde gözlenen en küçük değere yakın bir noktadan
başlasın. Grafiklerin okuyucuda oluşturduğu izlenimi karşılaştırın ve hangi
sunumun hangi durumda savunulabilir olduğunu tartışın.
\end{etkinlik}

\section{Bilimsel raporlama örneği}

``Öğrencilerin akademik kaygı puanlarının medyanı 53{,}5
($Q_1=50{,}5$, $Q_3=58{,}5$) olarak bulunmuştur. Dağılım büyük değerlere doğru
çarpıktır ve 85 puanı kutu grafiği kuralına göre aykırı değer adayıdır.
Ortalama $55{,}75$ ($SS=10{,}64$) olmakla birlikte bu değer uç gözleme
duyarlıdır. Veri kaydı doğrulandıktan sonra gözlem korunmuş ve sağlam özet
olarak medyan ile IQR öncelikle raporlanmıştır.''

Bu paragraf sayıların yanında dağılım biçimini, veri kontrolünü ve alınan
analiz kararını da görünür kılar.

\section{Literatür veri setiyle IBM SPSS laboratuvarı}
\index{Student Performance veri seti!betimsel istatistik}
\index{SPSS!Explore komutu}
\index{SPSS!Frequencies komutu}
\index{uç değer!SPSS incelemesi}

Bu uygulamada \emph{Student Performance} matematik dosyasındaki yıl sonu notu
\texttt{G3} ile devamsızlık sayısı \texttt{absences} betimlenecektir
\cite{cortez2008data,cortezsilva2008}. Aynı örneklemde yer alan bu iki nicel
değişken farklı dağılım özellikleri gösterdiğinden, merkez ve yayılım ölçüsü
seçimini karşılaştırmak için uygundur.

\begin{researchbox}
G3 puanı 0--20 aralığındaki yıl sonu notudur. Devamsızlık ise sıfırdan
başlayan bir sayım değişkenidir. İki değişken nicel olsa da aynı dağılım
biçimine sahip olmak veya aynı özetlerle raporlanmak zorunda değildir.
\end{researchbox}

\subsection{Araştırma görevi}

\begin{enumerate}
  \item Her değişken için geçerli ve eksik gözlem sayısını denetleyin.
  \item Ortalama, medyan, standart sapma, çeyrekler ve açıklığı karşılaştırın.
  \item Histogram ile kutu grafiğini sayısal özetlerle birlikte okuyun.
  \item Aşırı görünen değerleri otomatik olarak silmeden önce veri sözlüğünü inceleyin.
  \item G3 dağılımını okul gruplarına göre ayrı nokta veya kutu grafikleriyle karşılaştırın.
\end{enumerate}

\subsection{SPSS syntax}

Birinci bölümde içe aktarılan \texttt{StudentMath} veri seti açıkken aşağıdaki
syntax çalıştırılır. \textbf{Frequencies} kesin yer ve yayılım özetlerini,
\textbf{Explore} ise histogram, kutu grafiği, normal Q--Q grafiği ve en uç
gözlem listesini üretir.

\begin{lstlisting}[language={}]
DATASET ACTIVATE StudentMath.

FREQUENCIES VARIABLES=G3 absences
 /STATISTICS=MEAN MEDIAN STDDEV QUARTILES MINIMUM MAXIMUM
 /HISTOGRAM NORMAL
 /ORDER=ANALYSIS.

EXAMINE VARIABLES=G3 absences
 /PLOT BOXPLOT HISTOGRAM NPPLOT
 /STATISTICS DESCRIPTIVES EXTREME(5)
 /CINTERVAL 95
 /MISSING LISTWISE.

EXAMINE VARIABLES=G3 BY school
 /PLOT BOXPLOT
 /COMPARE GROUPS
 /STATISTICS DESCRIPTIVES
 /CINTERVAL 95
 /MISSING LISTWISE
 /NOTOTAL.
\end{lstlisting}

Menüyle çalışmak için \textbf{Analyze $>$ Descriptive Statistics $>$
Frequencies} yolunda ortalama, medyan, standart sapma ve çeyrekler;
\textbf{Charts} altında histogram seçilir. Daha ayrıntılı incelemede
\textbf{Analyze $>$ Descriptive Statistics $>$ Explore} penceresinde G3 ve
devamsızlık \textbf{Dependent List} alanına taşınır. Okula göre ayrı sonuç
istenirse \texttt{school}, \textbf{Factor List} alanına eklenir.

\subsection{SPSS çıktısının erişilebilir özeti}

\begin{table}[htbp]
\centering
\caption{G3 ve devamsızlık değişkenleri için SPSS betimsel özetleri.}
\begin{tabular}{@{}lrrrrrrrr@{}}
\toprule
Değişken & $n$ & Ort. & Md & SS & $Q_1$ & $Q_3$ & Min. & Maks. \\
\midrule
Yıl sonu notu (G3) & 395 & 10,42 & 11 & 4,58 & 8 & 14 & 0 & 20 \\
Devamsızlık & 395 & 5,71 & 4 & 8,00 & 0 & 8 & 0 & 75 \\
\bottomrule
\end{tabular}
\end{table}

G3 için çeyrekler arası açıklık
\[
IQR=14-8=6
\]
ve devamsızlık için
\[
IQR=8-0=8
\]
olarak hesaplanır. SPSS sürümü ve yüzdelik ara değerleme yöntemi belirtilirse
çeyreklerde görülebilecek küçük yazılım farkları doğru biçimde belgelenmiş
olur.

\subsection{Çıktının analizi}

G3 ortalaması 10,42, medyanı 11'dir. İki merkez ölçüsünün yakın olması tek
başına normalliği kanıtlamaz. Histogramda 0 puanındaki yığılma ve puanların
0--20 ile sınırlı olması ayrıca görülmelidir. Tukey kutu grafiği sınırları
yaklaşık $8-1{,}5(6)=-1$ ile $14+1{,}5(6)=23$ olduğundan 0 ve 20 puanları
yalnız bu kurala göre aykırı değer sayılmaz. Bununla birlikte 0 puanlarının
başarısızlık, sınava girmeme veya başka bir süreçle ilişkili olup olmadığı veri
sözlüğü ve araştırma bağlamıyla açıklanmalıdır.

Devamsızlıkta ortalamanın 5,71, medyanın 4 olması ve maksimumun 75'e ulaşması
sağa çarpık bir dağılıma işaret eder. Kutu grafiğinin üst inceleme sınırı
\[
Q_3+1{,}5IQR=8+1{,}5(8)=20
\]
olduğundan 20'nin üzerindeki değerler aykırı değer adayı olarak işaretlenebilir.
Bu değerler olanaksız değildir; öğrencinin gerçekten uzun süre devamsızlık
yapmış olması mümkündür. Bu nedenle silme yerine önce kayıt doğrulama ve
duyarlılık incelemesi gerekir.

G3 için ortalama ve standart sapma, medyan ve IQR ile birlikte verilebilir.
Devamsızlık için ise medyan ve IQR dağılımın tipik düzeyini daha sağlam biçimde
özetler; ortalama ve standart sapma karşılaştırma amacıyla ikincil olarak
sunulabilir. Böylece iki değişken için yazılımın aynı tabloyu üretmesi,
araştırmacının aynı raporlama kararını vermesini zorunlu kılmaz.

\begin{outputguidebox}
Explore çıktısını şu sırayla okuyun: (1) geçerli ve eksik $n$, (2) minimum ve
maksimumun veri sözlüğüyle uyumu, (3) ortalama--medyan karşılaştırması,
(4) standart sapma ve IQR, (5) histogram, (6) kutu grafiğindeki gözlem
numaraları, (7) Q--Q grafiği. Yalnız Shapiro--Wilk tablosuna bakarak betimsel
yorum yapılmamalıdır.
\end{outputguidebox}

\subsection{Örnek bulgu paragrafı}

``Matematik yıl sonu notlarının ortalaması 10,42 ($SS=4{,}58$), medyanı
11'dir ($Q_1=8$, $Q_3=14$, $n=395$). Puanlar 0--20 aralığındadır ve
histogramda 0 puanında yığılma görülmüştür. Devamsızlığın medyanı 4'tür
($Q_1=0$, $Q_3=8$); ortalama 5,71 ($SS=8{,}00$) ve aralık 0--75'tir.
Devamsızlık dağılımı sağa çarpık olduğundan merkez ve yayılım için medyan ile
IQR öncelikle yorumlanmış, 20'nin üzerindeki gözlemler silinmeden önce kayıt
ve bağlam açısından incelemeye alınmıştır.''

\section{R ile çözümlü uygulama}
\label{sec:r-b03}

\textbf{Soru.} Çalışma süreleri 2, 3, 3, 4 ve 13 saat olan örnekte merkez,
yayılım ve aykırı değer sınırlarını bulun. Bölümde kullanılan doğrusal
enterpolasyonla uyum için yüzdelik hesabında \texttt{type = 7} seçin.

\begin{lstlisting}[style=prismR]
saat <- c(2, 3, 3, 4, 13)
ceyrek <- quantile(saat, c(0.25, 0.75), type = 7)
ceyrek_acikligi <- IQR(saat, type = 7)
c(ortalama = mean(saat), medyan = median(saat),
  varyans = var(saat), standart_sapma = sd(saat))
quantile(saat, c(0, 0.25, 0.5, 0.75, 1), type = 7)
sinir <- c(ceyrek[1] - 1.5 * ceyrek_acikligi,
           ceyrek[2] + 1.5 * ceyrek_acikligi)
sinir
saat[saat < sinir[1] | saat > sinir[2]]
hist(saat, breaks = c(0, 3, 6, 9, 12, 15), right = FALSE,
     main = "Calisma suresi", xlab = "Saat")
boxplot(saat, horizontal = TRUE, xlab = "Saat")
c(ortalama_13_haric = mean(saat[saat != 13]),
  medyan_13_haric = median(saat[saat != 13]))
\end{lstlisting}

\begin{outputguidebox}
\textbf{Çözüm ve kontrol değerleri.} Ortalama 5, medyan 3, örneklem
varyansı 20,5 ve örneklem standart sapması 4,5277'dir. Beş sayı özeti
$(2;3;3;4;13)$; çeyrekler açıklığı 1; alt ve üst sınırlar 1,5 ve 5,5'tir.
13 saat üst sınırı aşar. Bu değer çıkarıldığında ortalama 3'e iner;
medyan 3 olarak kalır.
\end{outputguidebox}

\textbf{Yorum.} \texttt{sd} ve \texttt{var} örneklem hesabında $n-1$
bölenini kullanır. \texttt{boxplot} kutu sınırları için Tukey menteşelerini
kullanır; bunlar bu örnekte seçilen çeyreklerle aynıdır, her örneklemde aynı
olmak zorunda değildir. Aykırı değer işareti otomatik silme gerekçesi değildir.

\textbf{Pekiştirme ve yanıt.} 13'ü çıkarma işlemi yalnız duyarlılık
karşılaştırmasıdır. Ortalama iki saat değişirken medyanın değişmemesi,
medyanın bu uç değere karşı daha sağlam olduğunu gösterir.
\section{Bölüm sonu çözümlü problemler}

\begin{enumerate}
  \item \textbf{Merkez ve yayılma.} Beş öğrencinin oturum sayıları $2,3,3,4,8$'dir. Ortalama, medyan, ranj ve örneklem standart sapmasını hesaplayın.

  \textbf{Çözüm:} Ortalama $\bar{x}=20/5=4$, medyan 3 ve ranj $8-2=6$'dır. Kareli sapmalar toplamı $(-2)^2+(-1)^2+(-1)^2+0^2+4^2=22$ olduğundan örneklem varyansı $22/(5-1)=5{,}5$, standart sapma $s=\sqrt{5{,}5}\approx2{,}35$ olur.

  \item \textbf{Çeyrekler ve aykırı değer.} Sıralı puanlar $4,5,6,6,7,8,9,15$ olsun. Ortanca-yarılar yöntemiyle $Q_1$, $Q_3$, IQR ve üst aykırı değer sınırını bulun.

  \textbf{Çözüm:} Alt yarının medyanı $Q_1=(5+6)/2=5{,}5$, üst yarının medyanı $Q_3=(8+9)/2=8{,}5$'tir. $IQR=3$ ve üst sınır $8{,}5+1{,}5(3)=13$ olur. 15 aykırı değer adayıdır; otomatik olarak silinmez.

  \item \textbf{Sınıflandırılmış ortalama.} Puan sınıflarının orta noktaları 15, 25 ve 35; frekansları sırasıyla 4, 10 ve 6'dır. Yaklaşık ortalamayı bulun.

  \textbf{Çözüm:} Toplam frekans $n=20$ ve ağırlıklı toplam $4(15)+10(25)+6(35)=520$'dir. Sınıflandırılmış ortalama $\bar{x}_g=520/20=26$ olur. Bu değer, her sınıftaki gözlemlerin orta noktada bulunduğu varsayımına dayanan yaklaşık bir özettir.
\end{enumerate}

\bolumkaynaklari{\cite{anscombe1973,degroot2012,wasserman2004,cortez2008data,cortezsilva2008}}

\section*{Hatırla--Uygula--Yansıt}

\begin{enumerate}
\renewcommand{\labelenumi}{\textbf{\Alph{enumi}.}}
  \item \textbf{Hatırla:} Ortalama, medyan ve modun güçlü ve sınırlı yönlerini karşılaştırın.
  \item \textbf{Hesapla:} $4,5,5,6,10$ verisi için ortalama, medyan, açıklık ve örneklem standart sapmasını bulun.
  \item \textbf{Sınıflandırılmış veri:} Bir frekans tablosunda sınıf orta noktalarını kullanarak yaklaşık ortalama ve standart sapmanın nasıl hesaplandığını açıklayın; medyan sınıfını kümülatif frekanstan belirleyin.
  \item \textbf{Seç:} Sağa çarpık danışma bekleme süreleri için uygun merkez, yayılım ve grafik üçlüsünü gerekçelendirin.
  \item \textbf{Eleştir:} Y ekseni 48'den başlayan ve grup ortalamaları 50 ile 52 olan bir çubuk grafiğin oluşturabileceği yanıltıcı izlenimi açıklayın.
  \item \textbf{Uygula:} En az 15 gözlemden oluşan küçük bir veri setini Python veya SPSS ile özetleyip bir tablo ve bir grafik üretin.
  \item \textbf{Yansıt:} PDR verilerinde aykırı bir gözlemi silmenin bilimsel ve mesleki sonuçlarını tartışın.
\end{enumerate}